\scanpage{28}

technically \unsure{simple|flexible}. If $X$ is \deleted{\unsure{paracompact}} \inserted{a} paracompact space, one will be able to take (provisionally !) for $\mathcal V_{Y,X}$ the ``germ \unsure{of spaces|of space}'' of $X$ around $Y$ -- which is therefore a topos, namely the topos-theoretic $\varprojlim$ of the open neighborhoods of $Y$ in $X$. Intuitively, $\mathcal V_{Y,X}\hookrightarrow X$ has the properties of an open immersion, and $Y\hookrightarrow\unsure{\mathcal V_{Y,X}|\deleted{\mathcal V}_{Y,X}}$ those of a closed immersion, and (3) constitutes a canonical factorization of (1).

% The cancellation has a long top stroke, an open right side, a stepped lower-left edge and four diagonals; the words remain ordinary transcribed text.
\deleted{The hypothesis of normal local triviality implies that one has (under suitable conditions of ``tameness'') an inclusion}

\margin{The main interest of this factorization is that \unsure{the inclusion}''
\begin{tikzcd}[ampersand replacement=\&,column sep=small,cramped]
Y\arrow[r,-,to path={(\tikztostart)--(\tikztotarget) node[circle,fill=black,inner sep=.45pt] {} \tikztonodes}]\&\mathcal V_{Y,X}
\end{tikzcd}
is a homotopy equivalence'' (independently of any hypothesis of local triviality)}

(4)
\diagram{\begin{tikzcd}[ampersand replacement=\&]
Y\arrow[r,"\text{homotopy equivalence}"]\&\mathcal V_{Y,X}
\end{tikzcd}}

Suppose for example $Y$ and $X-Y=X^*$ \deleted{some} manifolds, and $Z\inserted{{}^*=Z\setminus\{z\}}$ a manifold (under hypotheses of tameness, the hypothesis on $Z$ should be a consequence of the first hypothesis), \deleted{\unsure{$Z^*$ is|and that}} whereas $z$ is a singular point \inserted{of $Z$}, i.e. $X$ is ``singular along $Y$''. It is natural then to want to \unsure{reconstruct|reconstruct} $X$ in terms of the manifolds $Y$ and $X^*$,

\margin{27}
